\section{Near a vertex}\label{sec:vertex}

This section is the worked case of the boundary layer of Section~\ref{sec:boundary-layer}
(Theorem~\ref{thm:layer} and Proposition~\ref{prop:pair-face}) on the complete graph near a
vertex, and Section~\ref{sec:layer-vertex} shows that with one rare count it is the discrete
pair-face layer. The vertex law covers the regime in which $T$ is of order $\sqrt N$, where the
lattice parameter is of order one and the Laguerre kernel cannot be replaced by the Bessel
kernel of the continuum layer.

We keep the notation of Section~\ref{sec:completegraph}, namely $Q=J-dI$, the uniform start,
$p$, $u$, $v$, and $S_n(u)=P_N(n)/V_N$ by~\eqref{eq:SnVN}. Let one count $b$ be large and let
the other positive counts $a_1,\ldots,a_\ell$ be small compared with $b$, with $\ell\ge1$ fixed. We
put $|a|=a_1+\cdots+a_\ell$ and write $S_{(b,a)}$ for $S_n$ with $n=(b,a_1,\ldots,a_\ell)$. These
rare counts are not the constants $a_k$ of Section~\ref{sec:completegraph}. Since $S_n$ does not
depend on the start, $P_N(n)\le(\max_i\mu_i)p^{N-1}S_n(u)$, so upper bounds on $S_n$ hold for
every start.

A rare run away from the ends needs two switches, which gives a factor $u^2$, and it can go in
about $b$ places. So the rare counts enter through $bu^2$, and for fixed rare counts the
crossing moves to odds $u$ of order $b^{-1/2}$, that is, to $T$ of order $\sqrt N$. Use $B_a$
from~\eqref{eq:refresh-polynomial} and put
\[
\mathcal I(x)=\sum_{j\ge0}\frac{x^j}{j!\,(j+1)!}=\frac{I_1(2\sqrt x)}{\sqrt x},
\]
where $I_1$ is the modified Bessel function~\cite[Eq.~10.25.2]{dlmf}, and $\mathcal I(0)=1$.
By the Laguerre form in~\eqref{eq:refresh-polynomial}, the next lemma is a quantitative form
of the limit of Laguerre polynomials to Bessel functions.

\begin{lemma}[a uniform Laguerre--Bessel inequality]\label{lem:uniform-laguerre-bessel}
For every integer $a\ge1$ and real $w>0$,
\begin{equation}\label{eq:uniform-laguerre-bessel}
0\le1-\frac{B_a(w)}{w\,\mathcal I(aw)}\le\frac w2 .
\end{equation}
\end{lemma}

The proof is in Appendix~\ref{app:vertex}. Put
\[
\mathcal B(y)=\prod_{i=1}^{\ell}B_{a_i}(y),\qquad
\kappa(y)=\frac{y\,\mathcal B'(y)}{\mathcal B(y)} .
\]
The polynomial $\mathcal B$ has positive coefficients, $\mathcal B(0)=0$ and
$\mathcal B(1)\ge1$, so $\mathcal B(y_a)=1$ has exactly one positive root $y_a$, and $y_a\le1$.
The polynomial $S_{(b,a)}$ has nonnegative coefficients, zero constant term and a positive
coefficient at $u^\ell$, so $S_{(b,a)}(u)=1$ has exactly one positive root $u_{b,a}$, the
crossing odds of the composition $(b,a_1,\ldots,a_\ell)$ with a vertex.

\begin{theorem}[the vertex law]\label{thm:vertex-uniform-root}
Fix $\ell\ge1$ and put
\[
 c_\ell=\ell+\sqrt\ell+1,\qquad
 \varepsilon=\sqrt{\frac{|a|+1}{b}},\qquad \delta=16c_\ell\varepsilon.
\]
For all positive counts $b,a_1,\ldots,a_\ell$ with $\delta\le1/2$,
\begin{equation}\label{eq:vertex-uniform-bracket}
 (1-\delta)\sqrt{y_a/b}\ \le u_{b,a}\le
 (1+\delta)\sqrt{y_a/b},\qquad \mathcal B(y_a)=1.
\end{equation}
In particular
\begin{equation}\label{eq:vertex-uniform-root}
 \biggl|\frac{u_{b,a}}{\sqrt{y_a/b}}-1\biggr|\le16c_\ell\sqrt{\frac{|a|+1}{b}}
 \qquad\text{whenever }\sqrt{\frac{|a|+1}{b}}\le\frac1{32c_\ell},
\end{equation}
so $u_{b,a}=\sqrt{y_a/b}\,(1+o(1))$ as $(|a|+1)/b\to0$. The constants are not
optimized.
\end{theorem}

For fixed rare counts Lemma~\ref{lem:vertex-randomization} gives
$S_{(b,a)}(\sqrt{y/b})\to\mathcal B(y)$ as $b\to\infty$, and the theorem
shows that the crossing stays near $\sqrt{y_a/b}$ uniformly as long as $|a|=o(b)$. For $d=2$
and $\ell=1$ its leading term is that of the fixed-edge expansion
of~\cite[Theorem~4.10]{paperA}. The proof is in Appendix~\ref{app:vertex}. It rests on two
facts. First, $\log\mathcal B$ is concave, and for $0<y\le1$ its elasticity $\kappa(y)$ lies
between $\max\{\ell,\frac12\sqrt{|a|y}\}$ and $\ell+\sqrt{\ell|a|y}$
(Lemma~\ref{lem:vertex-elasticity}). Second, $S_{(b,a)}(u)=(1-u)^{|a|}\E[\mathcal B(v\mathcal X)]$
exactly, where $\mathcal X$ is a gamma variable plus a compound binomial sum with mean
$2+(b-1)u$ (Lemma~\ref{lem:vertex-randomization}). Jensen's inequality then gives a lower bound
that exceeds $1$ at $u=(1+\delta)\sqrt{y_a/b}$, and the tangent line of $\log\mathcal B$ at
$bu^2$ with the moment generating function of $\mathcal X$ gives an upper bound that is below
$1$ at $u=(1-\delta)\sqrt{y_a/b}$.

\begin{remark}[the boundary of the simplex]\label{rem:B-boundary}
Away from a vertex, \Bversion~\cite{paperB} gives a Bessel majorant for every positive count
vector~\cite[Theorem~6.2]{paperB} and shows that, when at least two coordinates are of order
$N$, a small coordinate starts to matter at size of order $N/(\log N)^2$~\cite[Remark~6.3 and
Theorem~C.1]{paperB}. We do not use these results.
\end{remark}
